\section{An integral projective proof of adjoint positivity}
\label{p1logalt:section}

We give another proof of adjoint positivity for the whole highest Hodge
line of a rationally polarized integral pure variation.  The conclusion
is the same adjoint conclusion as Theorem~\ref{thm:adjoint} for the
rational parabolic extension itself. A positive rational multiple
has the same conclusion after rescaling the large adjoint parameter.
We first replace the
variation by an integral tensor construction, then descend it through
its full period map and subtract the boundary by comparing intersection
growth.

The tensor construction removes every monodromy element acting trivially
on the moving projective line.  This makes a transverse loop detect
strict loss of the line's numerical dimension at each remaining boundary
divisor.  The full period map supplies an algebraic base; the differential
of the map remembering only the highest line controls the boundary
subtraction.  These two ranks need not agree.

Let $P$ be a smooth connected projective complex variety, let $D$ be a reduced simple normal crossing
divisor, and put $P^\circ=P\setminus D$.  A \emph{top Hodge line} of a rationally
polarized integral pure variation $\mathbb V$ on $P^\circ$ means
\[
 F^p\mathbb V_{\cO}\quad\text{with}\quad
 F^{p+1}\mathbb V_{\cO}=0,\qquad
 \rank F^p\mathbb V_{\cO}=1.
\]
Here an integral variation has a local system of finite-rank free
$\mathbb Z$-modules, and the polarization is rational and flat.  The
weight and $p$ may be arbitrary integers.  The line itself is a complex
Hodge subbundle; no rationality assumption is made on it.  If $V$ is a flat reference
fiber, its \emph{line period map} is the equivariant holomorphic map
\[
 \ell:\widetilde{P^\circ}\longrightarrow\mathbb P(V),
 \qquad x\longmapsto F^p_x.
\]
Its generic differential rank does not depend on the reference fiber.

\begin{theorem}[Adjoint positivity]\label{p1logalt:positivity}
Let $P$ be a smooth connected projective complex variety and $D$ a
reduced simple normal crossing divisor.  Let $\mathbb V$ be a rationally
polarized integral pure variation on $P\setminus D$ whose highest
nonzero Hodge filtration piece is a line, and let
$L\in\operatorname{Pic}(P)\otimes\Q$ be its rational parabolic extension.  There are a projective birational
morphism $\pi:\widetilde P\to P$, with $\widetilde P$ smooth, a smooth
projective variety $Z$, a surjective morphism $h:\widetilde P\to Z$ with
connected fibers, and a nef rational divisor $A$ on $Z$ such that
\[
 \pi^*L\sim_{\Q}h^*A,
 \qquad K_Z+aA\ \text{is big for every sufficiently large rational }a.
\]
Here a divisor on a point is big.  Further smooth birational modifications
of the source are permitted by replacing $h$ with its composition with
the modification.
\end{theorem}

We first specify the extension and metric facts used in the proof.

\subsection{Extensions and numerical dimension}

The monodromy theorem makes the local monodromies quasi-unipotent.
For unipotent local monodromy, $L$ is the top filtration subbundle in the
Deligne--Schmid extension.  For quasi-unipotent monodromy, its rational
extension is characterized by the following rule: after an adapted cover
making local monodromy unipotent, its pullback is the unipotent extension,
as a rational line bundle.  This convention includes the parabolic
weights; it is not an arbitrary choice of a Deligne lattice.

\begin{lemma}[Extension facts]\label{p1logalt:extensions}
The rational extension has the following properties.
\begin{enumerate}[label=\textup{(\roman*)}]
\item It is nef and is compatible with positive tensor powers and
pullback by morphisms of smooth projective log pairs whose interiors map
to the original interior.
\item In the unipotent case, on each connected boundary stratum the local
limiting mixed Hodge structures determine a unique weight-graded piece
containing the top Hodge line.  These pieces glue to a polarizable pure
variation with rank-one top filtration, whose extended top line on the
closure of the stratum is the restriction of $L$.
\item The preceding restriction statement remains valid after an adapted
cover in the quasi-unipotent case.  In particular, numerical dimensions
of restrictions may be computed on a smooth adapted cover of the
stratum.
\end{enumerate}
\end{lemma}

\begin{proof}
In the unipotent case the canonical extension and relative monodromy
weight filtration are supplied by nilpotent-orbit theory
\cite{Schmid,CKS}; precise formulations for the top
line are given in \cite[\S 2.5.3, Lemma 2.16, Lemma 5.6]{BFMT}.  The
graded pieces are polarized by the primitive decomposition.

For a boundary stratum, take the weight filtration associated with an
interior point of its monodromy cone; it is independent of that choice.
Strictness and $\rank F^p=1$ select one nonzero highest graded piece.
The global object in (ii) is this pure graded variation.  To check its
gluing when the normal bundle is nontrivial, choose a local normal
coordinate $t$ and write $N$ for its monodromy logarithm.  Replacing $t$
by $t'=u(s)t$, with $u$ invertible along the stratum, changes the local
limiting filtration by
\[
 \exp\left(-\frac{\log u(s)}{2\pi\sqrt{-1}}N\right).
\]
Since $N$ lowers the monodromy weight filtration by two, this transition
acts identically on each weight graded.  For several normal coordinates
the same calculation uses the sum of the commuting monodromy logarithms,
which also acts trivially on their common weight graded.  The pure
graded data therefore glue.  Strictness identifies the rank-one top
filtration with its unique nonzero top graded piece.  At the remaining
corners, \cite[\S 2.5.3 and Lemma 2.16(3)]{BFMT} identifies the canonical
extensions on the entire smooth stratum closure, so the identification
of lines holds there as well.

For completeness, the rational convention is compatible with these
operations as follows.  In a crossing chart, take sufficiently divisible
roots of its boundary coordinates.  The resulting unipotent extension
is equivariant for the finite deck group.  A common positive power kills
the characters of the stabilizers on its line fibers and descends to a
line bundle in the original chart.  The descents agree on overlaps:
compare them on a common root cover and use the uniqueness of the
unipotent extension.  This defines a class in
$\operatorname{Pic}(P)\otimes\Q$ and proves compatibility with further
pullbacks and powers.  One can perform the comparison globally on a
smooth adapted alteration; auxiliary branch divisors cause no difficulty
because the original variation extends over them.  Nefness descends
under a proper surjective morphism, since every curve downstairs is
dominated by a curve upstairs.  Finally, nef numerical dimension is
unchanged by a generically finite pullback, so the same construction on
strata proves (iii).
\end{proof}

\begin{lemma}[Curvature and intersections]\label{p1logalt:rank}
The numerical dimension of the rational top Hodge line equals the
generic rank of its line period map:
\[
 \nu(L)=\rank(d\ell)_{\mathrm{gen}}.
\]
In the unipotent case, if $d=\dim P$, $\omega$ is a smooth K\"ahler form,
and $\alpha$ is the first Chern form of the Hodge metric on $L|_{P^\circ}$,
then, for $0\le k\le d$,
\begin{equation}\label{p1logalt:chern}
 \int_{P^\circ}\alpha^k\wedge\omega^{d-k}
 =c_1(L)^k\cdot[\omega]^{d-k}.
\end{equation}
The integrals are absolutely convergent.
\end{lemma}

\begin{proof}
The closed-test-form version of this calculation is proved in
Lemma~\ref{p1varhd:curvature}; that proof uses only the extension facts in
Lemma~\ref{p1logalt:extensions}, not any descent or rank-drop result.
Taking the test form to be $\omega^{d-k}$ gives
Equation~\eqref{p1logalt:chern}.  Griffiths' curvature formula gives
\begin{equation}\label{p1logalt:curvature-kernel}
                 \ker\alpha=\ker(d\ell).
\end{equation}
If $r$ is the generic differential rank, then $\alpha^{r+1}=0$ everywhere
on the open set, whereas $\alpha^r\wedge\omega^{d-r}$ is positive on a
nonempty open subset.  The intersection identity and nefness therefore
give $\nu(L)=r$.  For quasi-unipotent monodromy, both numerical dimension
and differential rank are unchanged by a smooth adapted generically
finite cover, and Lemma~\ref{p1logalt:extensions} identifies the pulled-back
rational line with its unipotent extension.
\end{proof}

\subsection{Eliminating the projective monodromy kernel}

For a variation on a connected smooth algebraic variety, we will use the
condition
\begin{equation}\label{p1logalt:faithful-line}
 \gamma\ell(x)=\ell(x)\text{ for every }x
 \quad\Longrightarrow\quad \gamma=\id
 \quad(\gamma\text{ in the monodromy image}).
\end{equation}
This is a condition on the image of the moving line, rather than on all
of the ambient projective space.

\begin{lemma}[An integral tensor replacement]\label{p1logalt:kernel}
Let $U$ be a smooth connected complex algebraic variety, and
let $\mathbb V$ be a rationally polarized integral pure variation on $U$
whose highest nonzero Hodge filtration piece is a line $L^\circ$.
There are an integer $b>0$ and a rationally polarized integral pure
variation $\mathbb V_1$ on $U$ satisfying
Condition~\eqref{p1logalt:faithful-line}, whose highest nonzero Hodge
filtration piece is $(L^\circ)^{\otimes b}$.  On any smooth projective
simple-normal-crossing compactification of $U$, the rational parabolic
highest line of $\mathbb V_1$ is the $b$th power of that of $\mathbb V$.
\end{lemma}

\begin{proof}
For a general algebraic $U$, choose a dense quasi-projective open
$U_0$.  The map $\pi_1(U_0)\to\pi_1(U)$ is surjective, since loops
can be moved off the complement.  Thus the monodromy decompositions
below are unchanged, and a flat projector of Hodge type $(0,0)$ on
$U_0$ has that type throughout $U$ by continuity.
Write $\Gamma$ for the original monodromy image.  By semisimplicity
\cite[Theorem~4.2.6 and the added-in-proof note]{DeligneHodgeII}, the complex local system is a sum of
monodromy-isotypic summands $V_j$. These are complex Hodge subvariations. To see this
precisely, the fixed-part theorem equips
$\operatorname{End}_{\Gamma}(V)$ with its Hodge algebra structure, and its
center has type $(0,0)$ by \cite[Corollary 4.2.8(i),(ii)]{DeligneHodgeII}.
Equivalently, the connected Hodge action fixes every primitive central
idempotent of this semisimple algebra.  The images of those idempotents
are exactly the $V_j$.

The top line belongs to one summand $V_{j_0}$.  Retain that summand and
all its rational conjugates.  Their sum is a rational subvariation:
the central idempotents are defined over a finite splitting field of the
rational commutant algebra, and the selected sum is Galois invariant.
It inherits the rational polarization and the lattice obtained by
intersecting with the original lattice.  Denote it by $V_0$.  None of
this construction requires the highest line, or its scalar character,
to be rational.

Suppose $\gamma\in\Gamma$ fixes the moving line projectively everywhere.
Its projective fixed locus is the finite union of the projectivizations
of its eigenspaces.  The connected line image lies in one of them, so
$\gamma$ is scalar on the linear span $S$ of that image.  Equivariance
makes $S$ a nonzero $\Gamma$-subrepresentation of $V_{j_0}$.  Write
$V_{j_0}=E\otimes M$, with $E$ irreducible and $\Gamma$ acting trivially
on $M$.  Since $S$ contains a copy of $E$, scalar action on $S$ implies
the same scalar action on $E$ and hence on all of $V_{j_0}$.  Rational
conjugation gives scalar action on every retained $V_j$.

It follows that $\gamma|_{V_0}$ preserves the Hodge filtration at any
point.  It also preserves the rational polarization and the lattice.
The stabilizer of a polarized Hodge structure in its real polarization
group is compact, since it preserves the positive definite Hodge metric.
Its intersection with the integral automorphism group is finite.
Therefore $\gamma|_{V_0}$ has finite order. For this particular
variation, one may take the exponent of that finite stabilizer at a
fixed base point. There is also a choice depending only on
$R=\rank V_0$.  Indeed, an eigenvalue of an integral finite-order matrix
of rank $R$ has order $q$ with $\varphi(q)\le R$, because its cyclotomic
minimal polynomial divides the characteristic polynomial.  The finite
set of such $q$ therefore gives the explicit permissible choice
\[
             b=\operatorname{lcm}\{q\geq1:\varphi(q)\leq R\}.
\]
Every scalar on every retained summand has $b$th power one.  This
uniform exponent is chosen before constructing the new variation.

Inside $V_0^{\otimes b}$ take
\[
 V_1=\bigoplus_{j\ \mathrm{retained}}V_j^{\otimes b}.
\]
The sum is taken inside $V_0^{\otimes b}$. The individual terms are
distinct direct summands in the tensor expansion, so the sum
is direct.  It is Galois invariant and hence rational, and each summand
is a Hodge subvariation.  Restricting the tensor polarization and
intersecting with the tensor lattice make it rationally polarized and
integral.  The summand
$V_{j_0}^{\otimes b}$ has top line $(L^\circ)^{\otimes b}$; all other
summands have strictly smaller highest Hodge index.  Thus this is the
unique top line of $V_1$.

Every element considered above acts trivially on $V_1$.  Conversely, if
an element of the new monodromy image fixes its moving top line, choose
a lift in $\Gamma$.  Injectivity of the Veronese map
$\mathbb P(V_{j_0})\to\mathbb P(V_{j_0}^{\otimes b})$ shows that the lift
fixes the original moving line.  It is therefore killed on $V_1$.
This proves Condition~\eqref{p1logalt:faithful-line}.  Compatibility of extensions
with tensor powers finishes the proof.
\end{proof}

We now descend the variation itself, using faithfulness to kill
monodromy along connected period fibres. Consequently we may
normalize the period image in the relative algebraic closure of its
function field in $\C(P)$ before estimating boundary ranks. The
finite-inertia contradiction in Section~\ref{sec:rankloss} required
the actual image; here the tensor replacement makes fibre monodromy
trivial once it fixes the moving line on an open set.

\begin{lemma}[Full-period descent]\label{p1logalt:descent}
After applying Lemma~\ref{p1logalt:kernel} and making smooth birational
modifications of the projective models, the resulting variation
descends over a dense open subset of a smooth projective variety $Z$.
Its descended full period map is generically immersive and satisfies
Condition~\eqref{p1logalt:faithful-line}.  If $A_1$ is its rational top Hodge
extension, then the original extension satisfies
\[
 b\pi^*L\sim_{\Q}h^*A_1
\]
for a morphism $h:\widetilde P\to Z$ with connected fibers.
\end{lemma}

\begin{proof}
Let $\mathcal D$ be the period domain for the chosen polarized lattice,
and choose an arithmetic group $\Gamma_{\mathrm{ar}}$ containing the
actual monodromy image $\Gamma$.  By \cite[Theorem 1.1]{BakkerBrunebarbeTsimerman}, the full
period map factors through an algebraic dominant morphism
$P^\circ\to Y_0$, where $Y_0$ is an irreducible quasi-projective variety
and $Y_0^{\mathrm{an}}\to\Gamma_{\mathrm{ar}}\backslash\mathcal D$ is
a closed analytic immersion.  Normalize $Y_0$ in the relative algebraic
closure of $\C(Y_0)$ in $\C(P)$.  This extension is finite, so the
normalization is finite over $Y_0$.  Its function field is relatively
algebraically closed in $\C(P)$; in characteristic zero this makes the
generic fiber geometrically connected.  Shrink
source and target so that the latter is smooth and the map is a smooth
topologically locally trivial fibration with connected fibers.  To
obtain this shrinking for the possibly nonproper source, stratify a
proper algebraic compactification of the map together with its source
boundary and restrict over the open base stratum.

A local lift of the full period map along a fiber takes values in one
$\Gamma_{\mathrm{ar}}$-orbit.  These orbits are discrete, including when
the quotient has finite stabilizers, so that lift is locally constant
in fiber directions.  Continuation around a fiber loop therefore fixes
the lifted full period.  Transport the loop to nearby fibers using the
local trivialization.  In flat coordinates its actual monodromy fixes
the top line on a nonempty analytic open subset of the source.  This
identity extends to its connected universal cover, and
Lemma~\ref{p1logalt:kernel} kills that monodromy.  The homotopy sequence of
the fibration now factors the integral representation through the
ordinary fundamental group of the base.  The filtration descends as well: on a
simply connected base chart the pulled-back flat bundle is trivial,
the filtration is constant on each connected fiber, and holomorphic
local sections of the smooth map show that the descended filtration is
holomorphic.  Its transversality and polarization follow by pullback.

The descended period map has the dimension of the original full period
image and is generically immersive.  Base loops lift to the source, so
the same argument and Condition~\eqref{p1logalt:faithful-line} show that the descended
variation retains that condition.  Shrinking introduces no change to
it: any monodromy element fixing the line on the smaller open fixes it
identically by analytic continuation.

Choose a smooth projective compactification of the base, resolve its
boundary, and resolve the rational map from $P$ to this compactification.
The resulting morphism $h:\widetilde P\to Z$ has connected generic
fiber and hence connected fibers, by Stein factorization and normality
of $Z$.  Enlarge the source boundary to contain the inverse image of
the chosen base boundary and resolve it.  The rational extension of
the original variation is unchanged where it already extended.
Lemma~\ref{p1logalt:extensions} identifies its $b$-th power with $h^*A_1$.
\end{proof}

\subsection{A strict rank drop at nontrivial monodromy}

\begin{lemma}[A transverse-loop criterion]\label{p4int:transverse-loop}
Let $U=\Delta\times\Delta^{d-1}$, with coordinates $(z,s)$, and let
$U^*=\Delta^*\times\Delta^{d-1}$.  Suppose that a holomorphic map
$\ell:\widetilde{U^*}\to\mathbb P(V)$ is equivariant for a linear
transformation $T\in\operatorname{GL}(V)$ under one positive turn around
$z=0$, and that $\rank(d\ell)\leq r$.  Suppose that a holomorphic map
$q:U\to\mathbb C^r$ has tangential differential of rank $r$ at a point
of $z=0$, and that on $U^*$ its lift satisfies
$\ker(d\ell)\subseteq\ker(dq)$.
Then $T$ fixes every line in the image of $\ell$ on a nonempty open
subset.  If $\ell$ is the restriction of a line period map on a connected
universal cover, $T$ fixes that entire line image.  This includes $r=0$,
when $q$ is the map to a point.
\end{lemma}

\begin{proof}
After shrinking around the given point, the submersion theorem provides
coordinates
\[
            (z,q_1,\ldots,q_r,u_1,\ldots,u_{d-r-1}).
\]
The coordinate $z$ is retained because $dq$ already has rank $r$ on its
zero divisor.  On the punctured neighborhood, the kernel inclusion and
the rank bound imply
$\rank(d\ell)=r$ and $\ker(dq)=\ker(d\ell)$.
Fix $q$ and $u$ and make one positive turn in the $z$ coordinate.
Along every lift of that loop, the derivative of $\ell$ vanishes, so its
value is constant.  Its endpoints are related by $T$, which consequently
fixes the initial line.  Varying the fixed coordinates proves the claim
on an open subset.  The projective equality $T\ell=\ell$ then extends by
the identity theorem to the connected universal cover.  For $r=0$, the
rank bound makes $\ell$ locally constant and the same loop argument
applies with no $q$ coordinates.
\end{proof}

\begin{lemma}[Boundary rank drop]\label{p1logalt:boundary-drop}
Let $Z$ be smooth projective, let $D$ be a reduced simple normal crossing
divisor, and let a polarizable integral pure variation on $Z\setminus D$
have a top Hodge line satisfying Condition~\eqref{p1logalt:faithful-line}.  Write $A_1$
for its rational extension and $r=\nu(A_1)$.  If $D_i$ is a component
whose transverse monodromy is not the identity, then
\[
 \nu(A_1|_{D_i})<r.
\]
\end{lemma}

\begin{proof}
If the boundary rank did not drop, suitable single-valued functions of
the moving line would have the same differential kernel as that line.
A transverse loop in one of their fibers would then force the monodromy
to fix the moving line on an open set, contradicting
Condition~\eqref{p1logalt:faithful-line} after analytic continuation.

Work at a general point of $D_i$, with coordinates $(z,s)$ on
$\Delta\times\Delta^{d-1}$ and $D_i=\{z=0\}$.  Write
$T=T_sT_u$ for the Jordan decomposition of the transverse monodromy and
put $N=\log T_u$.  Choose $e>0$ with $T_s^e=1$ and set $z=t^e$.
In flat coordinates the nilpotent orbit extension supplies a holomorphic
line map $\Psi$ on the root polydisk such that
\begin{equation}\label{p1logalt:untwist}
 \ell(t,s)=
 \exp\left(\frac{e\log t}{2\pi\sqrt{-1}}N\right)\Psi(t,s),
 \qquad
 \Psi(\zeta t,s)=T_s\Psi(t,s),
 \quad \zeta=\exp(2\pi\sqrt{-1}/e).
\end{equation}
These are equalities of lines; $\ell$ is evaluated on the logarithmic
cover when necessary.

Let $W=W(N)$, centered at the weight of the variation.  On the boundary
there is a unique index $k$ in which the top line has a nonzero
weight-graded image.  The map $N$ is a strict morphism of limiting mixed
Hodge structures of type $(-1,-1)$.  Since $F^{p+1}=0$, it gives
\begin{equation}\label{p1logalt:survives}
 F^p\cap NV=N(F^{p+1})=0.
\end{equation}
Thus the limiting line survives in $C=V/NV$.

We also need to preserve its derivative rank.  The primitive decomposition
of the fixed graded local system identifies
\[
 \operatorname{Gr}^W_k V
 =P_k\oplus N\operatorname{Gr}^W_{k+2}V,
 \qquad
 P_k\xrightarrow{\ \sim\ }\operatorname{Gr}^W_k C.
\]
Here $k$ is at least the original weight: lower graded pieces are
entirely in the image of $N$ and cannot contain the top line.  The second
summand has zero $F^p$, so the graded moving top line lies in the flat
primitive summand $P_k$.  Its projective derivative rank is therefore
unchanged by the displayed isomorphism.  Projecting the boundary line
in $C$ further to its relevant graded piece recovers precisely this pure
line period map.  By Lemmas~\ref{p1logalt:extensions} and~\ref{p1logalt:rank},
its rank is
\[
 r_i=\nu(A_1|_{D_i}).
\]
This equality may be checked on a global adapted cover: restriction of
the pulled-back line to a component dominating $D_i$ is its generically
finite pullback, which preserves numerical dimension.

Equation~\eqref{p1logalt:untwist} makes $\Psi(0,s)$ an eigenline of $T_s$.
Its eigenvalue $\lambda$ is constant on a sufficiently small connected
boundary neighborhood.  Let $E=C_\lambda$ and let $q_0:V\to E$ be the
quotient followed by the eigenspace projection.  The projective map
\[
 Q(t,s)=\mathbb P(q_0)(\Psi(t,s))
\]
is defined near the chosen boundary point and has boundary rank at
least $r_i$.  Since $N$ vanishes on $C$, it also equals
$\mathbb P(q_0)(\ell(t,s))$.  Moreover $T_s$ acts by a scalar on $E$,
so $Q(\zeta t,s)=Q(t,s)$.  Hence $Q$ descends holomorphically to the
original $(z,s)$-polydisk, including $z=0$.  This descent uses a
projective map; a choice of a fractional scalar frame is unnecessary.

Suppose $r_i\ge r$.  Choose $r$ affine coordinates of $Q$, denoted
$q_1,\ldots,q_r$, whose differentials are independent along $D_i$.
The map $q=(q_1,\ldots,q_r)$ factors locally through $\ell$ on the
punctured chart.  Its differential therefore satisfies the kernel
inclusion in Lemma~\ref{p4int:transverse-loop}.  The rank bound for
$\ell$ is $r$ by Lemma~\ref{p1logalt:rank}.  The transverse-loop lemma
shows that $T$ fixes the moving highest line everywhere on the connected
universal cover.  Condition~\eqref{p1logalt:faithful-line} gives
$T=\id$, contrary to the choice of the divisor.  If $r=0$, the same
lemma applies with no coordinate functions.  If $r>d-1$, the supposed
inequality $r_i\geq r$ is already impossible.
\end{proof}

\subsection{Removing the boundary from the adjoint class}

\begin{lemma}[A boundary subtraction estimate]\label{p1logalt:volume}
Let $Z$ be smooth projective of positive dimension $d$, let $A_1$ be a
nef rational divisor, and let $J=\sum_iD_i$ be reduced with smooth
components.  Suppose $K_Z+J$ is big and
\[
 \nu(A_1|_{D_i})<r:=\nu(A_1)\quad\text{for every }i.
\]
Then $K_Z+aA_1$ is big for every sufficiently large rational $a$.
\end{lemma}

\begin{proof}
If $r=0$, the restriction hypothesis forces $J=0$, so $K_Z$ is already
big.  Suppose $r>0$.  By Kodaira's lemma write
\[
 K_Z+J\sim_{\Q}H+E,\qquad H\text{ ample},\quad E\ge0,
\]
and set $B_a=H+aA_1$ for rational $a\ge0$.  Choose one ample integral
divisor $C_0$ such that both $C_0$ and $C_0+J$ are globally generated,
and put $H'=H+C_0$.  This choice is independent of $a$ and of all
section degrees below.

For every sufficiently divisible $m>0$, filter by $kJ$ for $0\le k<m$.
The restriction sequences give
\[
 h^0\bigl(Z,m(B_a-J)\bigr)
 \ge h^0(Z,mB_a)
 -\sum_{k=0}^{m-1}h^0\bigl(J,(mB_a-kJ)|_J\bigr).
\]
Since $J$ is reduced, its structure sheaf injects into
$\bigoplus_i\cO_{D_i}$.  Also
\[
 mC_0+kJ=(m-k)C_0+k(C_0+J)
\]
is globally generated.  On each $D_i$, multiplication by a section
which is not identically zero bounds the restricted summand by
$h^0(D_i,m(H'+aA_1)|_{D_i})$.  We obtain the exact estimate
\begin{equation}\label{p1logalt:uniform-sections}
 h^0\bigl(Z,m(B_a-J)\bigr)
 \ge h^0(Z,mB_a)
 -m\sum_i h^0\bigl(D_i,m(H'+aA_1)|_{D_i}\bigr).
\end{equation}
Fix $a$ and divide by $m^d/d!$.  Asymptotic Riemann--Roch for the ample
divisors involved, along divisible $m\to\infty$, yields
\begin{equation}\label{p1logalt:uniform-volume}
 \vol(B_a-J)
 \ge B_a^d-d\sum_i(H'+aA_1)^{d-1}\cdot D_i.
\end{equation}
No uniform Riemann--Roch remainder in $a$ is needed: the exact
inequality in Equation~\eqref{p1logalt:uniform-sections} holds for every $a$, and the
limit is taken with $a$ fixed.

The first polynomial on the right of
Equation~\eqref{p1logalt:uniform-volume} has degree $r$, with positive leading
coefficient $\binom dr A_1^rH^{d-r}$.  Every restriction polynomial has
degree at most $\nu(A_1|_{D_i})\le r-1$.  The right side is therefore
positive for every sufficiently large $a$.  Finally
$K_Z+aA_1\sim_{\Q}(B_a-J)+E$ is big.
\end{proof}

\begin{proof}[Proof of Theorem~\ref{p1logalt:positivity}]
Apply Lemmas~\ref{p1logalt:kernel} and~\ref{p1logalt:descent}.  We obtain a
variation on a dense open of a smooth projective $Z$ satisfying
Condition~\eqref{p1logalt:faithful-line}, with generically immersive full period map,
and rational Hodge extension $A_1$ satisfying
$b\pi^*L\sim_{\Q}h^*A_1$.  Resolve the boundary so that it is simple normal crossing.  Pullback
compatibility retains the line identity.  Delete every boundary component with identity transverse
monodromy, and call the remaining reduced divisor $J$.

The variation extends to $Z\setminus J$.  Indeed, at a point lying only
on deleted components all local monodromies are the identity, and the
nilpotent orbit extension with all monodromy logarithms zero is a pure
polarized extension.  The full period map is still generically immersive.
Thus \cite[Theorem~1.1 and Remark~1.2(1)]{BrunebarbeCadorel} applies
to the polarized variation on $Z\setminus J$ and gives the bigness of
$K_Z+J$; its hypothesis is immersivity at one point, and it imposes no
extra condition on boundary monodromy.
Lemma~\ref{p1logalt:boundary-drop} applies to every component of $J$, and
Lemma~\ref{p1logalt:volume} proves that $K_Z+aA_1$ is big for all sufficiently
large rational $a$.  When $\dim Z=0$ the assertion holds by convention.
Set $A=A_1/b$.  This is nef, satisfies
$\pi^*L\sim_{\Q}h^*A$, and has the same asserted adjoint positivity
after rescaling the large parameter.
\end{proof}

\begin{remark}\label{p1logalt:no-abundance}
The construction permits $\nu(A)<\dim Z$.  In particular, the algebraic
map $h$ is obtained from the full period map of an integral tensor
replacement; its fibers are not asserted to be the leaves of the line
period map.  Neither abundance nor semiampleness of $A$ occurs in the
argument.
\end{remark}
